% Generated draft: requires mathematical and source review.
\documentclass[11pt]{article}
\usepackage[margin=25mm]{geometry}
\usepackage{fontspec}
\setmainfont{Latin Modern Roman}
\usepackage{amsmath,amssymb,mathtools,unicode-math}
\setmathfont{Latin Modern Math}
\usepackage{enumitem,needspace,xurl,hyperref}
\setlength{\parindent}{0pt}
\setlength{\parskip}{5pt}
\setlength{\emergencystretch}{3em}
\newcommand{\N}{\mathbb{N}}
\newcommand{\Z}{\mathbb{Z}}
\newcommand{\Q}{\mathbb{Q}}
\newcommand{\R}{\mathbb{R}}
\newcommand{\C}{\mathbb{C}}
\newcommand{\F}{\mathbb{F}}
\newcommand{\A}{\mathbb{A}}
\newcommand{\PP}{\mathbb{P}}
\newcommand{\E}{\mathbb{E}}
\newcommand{\eps}{\varepsilon}
\newcommand{\dd}{\,\mathrm{d}}
\begin{document}
\section*{UnsolvedMath ID: 8800072}
\textbf{Source problem code:} catalogue entry
\subsection{Definitions and mathematical statement}
A digital-signature scheme consists of key generation $\operatorname{KeyGen}(1^\lambda)$, signing $\operatorname{Sign}_{sk}(m)$, and verification $\operatorname{Verify}_{pk}(m,\sigma)$. Correctness requires honestly generated signatures to verify, up to any explicitly allowed negligible error. A standard security target is existential unforgeability under adaptive chosen-message attack: an efficient adversary with a signing oracle produces, with only negligible probability, a valid signature on a message not queried to that oracle. Negligible means eventually smaller than $\lambda^{-c}$ for every fixed $c>0$.

Construct efficient signature schemes whose security reduces to explicitly specified ideal-lattice problems. An ideal lattice is the geometric image of an ideal in a specified number-field order, with its embedding and norm fixed; a candidate hardness problem and approximation parameters must be named. State the reduction, its computational/adversarial model, and key/signature sizes and runtime. Polynomial asymptotic efficiency and concrete practical efficiency are distinct and the source supplies no numerical benchmark.

The supplied summary reports successful ring/ideal-lattice constructions by Lyubashevsky and others and cites Dilithium as a practical standardized example. Its exact assumptions, reduction models, and publication identifiers are incomplete. In particular, module-lattice security is not automatically a rank-one ideal-lattice security statement, so the Dilithium mention does not by itself settle every narrow interpretation of the requested ideal-lattice assumption. No unconditional security or unprovided quantum/standard-model guarantee is added.
\subsection{Short English statement}
Build efficient signatures with a precise security reduction to specified ideal-lattice problems. Fix the security and oracle models and distinguish ideal-, ring-, and module-lattice assumptions when interpreting the supplied construction claims.
\subsection{Sources}
\begin{itemize}
\item Daniele Micciancio, IPAM \emph{Workshop on Number Theory and Cryptography --- Open Problems} (2006), attachment 6240, PDF page 38, schedule source \url{https://www.ipam.ucla.edu/programs/workshops/workshop-i-number-theory-and-cryptography-open-problems/?tab=schedule}; accessed 2026-07-29. The supplied provenance records a PDF/legacy-PPT visual audit.
\item Supplied construction attributions: Lyubashevsky (2012); Stehl\'e--Steinfeld--Tanaka--Xagawa, identified by the source summary as GST14; and Dilithium variants. No complete publication citations, exact assumptions, or security-model statements are supplied.
\item Local registry \texttt{data/unsolvedmath/problems.json}, permanent UnsolvedMath ID $8800072$, source code catalogue entry, \emph{Ideal-lattice digital signatures}. Dataset attribution: Ulam AI and the UnsolvedMath Contributors.
\end{itemize}
\end{document}
